% Book 4, Year 2, chapter 29 -- Learning from examples: decision trees and clustering.
% Provenance (never printed): sources/book4/BRIEFS.md, chapter 29.
\chapter{Learning from Examples: Decision Trees and Clustering}\label{ch:b2:trees-and-clustering}

The library would like to know, when a book goes out, whether it will come back late. Nobody has written down the rule,
if there is one; there are only past loans, each with a few facts and the answer. A program can look for the rule in
them:

\omcode{code/bachelor-2/29-trees-and-clustering/listings/id3.py}{74}{78}{Learn a \omterm{def:b2:trees-and-clustering:decision-tree}{decision tree} from 700 past loans and
test it on 300 others}\label{lst:b2:trees-and-clustering:hook}

\omoutput{code/bachelor-2/29-trees-and-clustering/outputs/id3.txt}{The tree the program found, and its score on loans it
never saw.}\label{out:b2:trees-and-clustering:hook}

The program was given five facts per loan: the days lent, the book's pages, the member's age, how many books the member
has already returned late, and how far from the library the member lives. It kept two of them and found a rule of two
questions: a member with at least two late returns who lives more than about 2.5 km away brings the book back late. The
loans were generated, and that is exactly the rule that was planted in them, so the program can be checked. Real data
are not so kind: labels are noisy, and a program that fits its \omterm{def:b2:trees-and-clustering:examples}{examples} too closely learns the noise. This chapter
builds the vocabulary and the tools to measure that, then turns to \omterm{def:b2:trees-and-clustering:examples}{examples} with no answer at all, where the task is to
find groups.

\begin{recall}
The grade 11 volume: the $k$-nearest-neighbour \omterm{def:b2:trees-and-clustering:learning}{classifier}, distances between points, training and test \omterm{def:b2:trees-and-clustering:examples}{examples}. The
Year 1 volume: trees and their traversals, Kruskal's algorithm for the minimum spanning tree, with union--find. The
mathematics books: probability of an event, mean, the logarithm in base 2. Chapter 27: seeded pseudo-random numbers.
\end{recall}

\section{Learning as an algorithmic problem}\label{sec:b2:trees-and-clustering:learning}

\begin{definition}[Example, feature, class label, training set, test set]\label{def:b2:trees-and-clustering:examples}
An \emph{example}\index{example (learning)} is a vector of values, its \emph{features}\index{feature}, possibly with a
\emph{class label}\index{class label}, the answer to be predicted, taken from a finite set. A \emph{training
set}\index{training set} is a set of labelled examples given to a learning algorithm; a \emph{test set}\index{test set}
is a set of labelled examples, drawn from the same source, that the algorithm never sees and that serves only to measure
the result.
\end{definition}

\begin{definition}[Supervised and unsupervised learning, classifier]\label{def:b2:trees-and-clustering:learning}
In \emph{supervised learning}\index{supervised learning}, an algorithm receives labelled \omterm{def:b2:trees-and-clustering:examples}{examples} and returns a
\emph{classifier}\index{classifier}, a function from \omterm{def:b2:trees-and-clustering:examples}{feature} vectors to \omterm{def:b2:trees-and-clustering:examples}{class labels}. In \emph{unsupervised
learning}\index{unsupervised learning}, the \omterm{def:b2:trees-and-clustering:examples}{examples} have no labels and the algorithm looks for structure in them, such
as groups.
\end{definition}

A \omterm{def:b2:trees-and-clustering:learning}{classifier} is useful only on \omterm{def:b2:trees-and-clustering:examples}{examples} it has not seen: a table of the training \omterm{def:b2:trees-and-clustering:examples}{examples} would answer all of them
correctly and say nothing about the next loan. Hence the \omterm{def:b2:trees-and-clustering:examples}{test set}, kept apart until the end.

\begin{definition}[Confusion matrix, accuracy]\label{def:b2:trees-and-clustering:confusion}
The \emph{confusion matrix}\index{confusion matrix} of a \omterm{def:b2:trees-and-clustering:learning}{classifier} on a labelled set counts, for each pair of classes
$(a, b)$, the \omterm{def:b2:trees-and-clustering:examples}{examples} of true class $a$ that it labels $b$. Its \emph{accuracy}\index{accuracy} is the fraction of the
\omterm{def:b2:trees-and-clustering:examples}{examples} on the diagonal, those labelled correctly.
\end{definition}

\omcode{code/bachelor-2/29-trees-and-clustering/listings/confusion.py}{5}{16}{The \omterm{def:b2:trees-and-clustering:confusion}{confusion matrix} of a tree of depth 2
on 300 test loans whose labels are 15~\% noise}\label{lst:b2:trees-and-clustering:confusion}

\omoutput{code/bachelor-2/29-trees-and-clustering/outputs/confusion.txt}{\omterm{def:b2:trees-and-clustering:confusion}{Accuracy} is one number; the matrix says which
errors are made.}\label{out:b2:trees-and-clustering:confusion}

\omterm{def:b2:trees-and-clustering:confusion}{Accuracy} alone can mislead when one class is rare: a \omterm{def:b2:trees-and-clustering:learning}{classifier} that always answers ``on time'' scores 0.777 here,
without learning anything. The matrix shows that the tree catches about half the late loans and raises few false alarms.
The late loans it misses are mostly noise: their label was flipped at random, and no rule can predict them.

\section{Nearest neighbours and $k$-d trees}\label{sec:b2:trees-and-clustering:knn}

\begin{definition}[$k$-nearest-neighbour classifier]\label{def:b2:trees-and-clustering:knn}
The \emph{$k$-nearest-neighbour classifier}\index{k-nearest-neighbour classifier@$k$-nearest-neighbour classifier}
labels a \omterm{def:b2:trees-and-clustering:examples}{feature} vector $q$ with the most frequent label among the $k$ training \omterm{def:b2:trees-and-clustering:examples}{examples} closest to $q$ in Euclidean
distance, breaking ties by the label of the nearest of them.
\end{definition}

The grade 11 volume met this \omterm{def:b2:trees-and-clustering:learning}{classifier} on two \omterm{def:b2:trees-and-clustering:examples}{features}; here it runs on five. A distance adds up differences of
\omterm{def:b2:trees-and-clustering:examples}{features} measured in different units, days, pages and metres, so the \omterm{def:b2:trees-and-clustering:examples}{features} are first rescaled to $[0, 1]$ by their
range on the \omterm{def:b2:trees-and-clustering:examples}{training set}; without it, the distance in metres would decide alone (\cref{exo:b2:trees-and-clustering:2}).

\omcode{code/bachelor-2/29-trees-and-clustering/listings/knn.py}{8}{30}{The \omterm{def:b2:trees-and-clustering:learning}{classifier}, scaling and evaluation}\label{lst:b2:trees-and-clustering:knn}

\omoutput{code/bachelor-2/29-trees-and-clustering/outputs/knn.txt}{\omterm{def:b2:trees-and-clustering:confusion}{Accuracy} by $k$, with all five \omterm{def:b2:trees-and-clustering:examples}{features} and with the
two that matter.}\label{out:b2:trees-and-clustering:knn}

With $k = 1$ the \omterm{def:b2:trees-and-clustering:learning}{classifier} copies the noise of the nearest \omterm{def:b2:trees-and-clustering:examples}{example}; a larger $k$ averages it out. The three \omterm{def:b2:trees-and-clustering:examples}{features}
that play no part in the rule still enter every distance and blur it: dropping them helps. A \omterm{def:b2:trees-and-clustering:decision-tree}{decision tree}, below,
chooses its \omterm{def:b2:trees-and-clustering:examples}{features} itself.

\begin{proposition}[Cost of $k$-nearest neighbours]\label{prop:b2:trees-and-clustering:knn-cost}
With $n$ training \omterm{def:b2:trees-and-clustering:examples}{examples} of $d$ \omterm{def:b2:trees-and-clustering:examples}{features}, the \omterm{def:b2:trees-and-clustering:learning}{classifier} needs no training beyond storing them, in $O(nd)$ memory, and
labelling one query by computing all distances costs $O(nd + n \log k)$ time.
\end{proposition}

\begin{proof}
Each of the $n$ distances costs $O(d)$. The $k$ smallest of $n$ numbers are found by a pass that keeps a heap of the $k$
best so far, at $O(\log k)$ per number; the vote then costs $O(k)$. (The listing sorts all $n$, at $O(n \log n)$.)
\end{proof}

Computing every distance is wasteful when the points live in a space of few dimensions: most are obviously far away.

\begin{definition}[$k$-d tree]\label{def:b2:trees-and-clustering:kd-tree}
A \emph{$k$-d tree}\index{k-d tree@$k$-d tree} of a set of points in $\mathbb{R}^k$ is a binary tree with one point per
node. A node at depth $i$ splits on coordinate $i \bmod k$: the points of its left subtree have that coordinate at most
the node's, those of its right subtree at least the node's. Built by medians, it takes the point of median coordinate as
the node, and so has depth $\lfloor \log_2 n \rfloor$.
\end{definition}

\omcode{code/bachelor-2/29-trees-and-clustering/listings/kdtree.py}{6}{35}{A $k$-d tree in the plane, and the search for
the nearest point}\label{lst:b2:trees-and-clustering:kdtree}

\begin{omfigure}
\begin{tikzpicture}[scale=0.55, split0/.style={omDef, thick}, split1/.style={omThm, thick},
  root/.style={circle, draw, thick, minimum size=9pt, inner sep=0pt}]
\draw[black!50] (0, 0) rectangle (10, 10);
% written by figs.py; do not edit
\draw[split0] (5.3, 0) -- (5.3, 10);
\draw[split1] (0, 2.0) -- (5.3, 2.0);
\draw[split0] (2.5, 0) -- (2.5, 2.0);
\draw[split1] (0, 1.3) -- (2.5, 1.3);
\draw[split1] (2.5, 1.8) -- (5.3, 1.8);
\draw[split0] (2.3, 2.0) -- (2.3, 10);
\draw[split1] (0, 4.1) -- (2.3, 4.1);
\draw[split1] (2.3, 6.2) -- (5.3, 6.2);
\draw[split1] (5.3, 6.6) -- (10, 6.6);
\draw[split0] (7.8, 0) -- (7.8, 6.6);
\draw[split1] (5.3, 2.4) -- (7.8, 2.4);
\draw[split1] (7.8, 1.1) -- (10, 1.1);
\draw[split0] (8.1, 6.6) -- (8.1, 10);
\draw[split1] (5.3, 8.3) -- (8.1, 8.3);
\draw[split1] (8.1, 7.8) -- (10, 7.8);
\fill (2.5, 1.3) circle (3.2pt);
\fill (4.0, 1.8) circle (3.2pt);
\fill (0.9, 4.1) circle (3.2pt);
\fill (8.9, 7.8) circle (3.2pt);
\fill (7.5, 2.4) circle (3.2pt);
\fill (5.3, 2.9) circle (3.2pt);
\fill (1.9, 1.3) circle (3.2pt);
\fill (2.3, 9.0) circle (3.2pt);
\fill (8.1, 7.9) circle (3.2pt);
\fill (7.8, 2.1) circle (3.2pt);
\fill (3.2, 6.2) circle (3.2pt);
\fill (7.2, 8.3) circle (3.2pt);
\fill (8.6, 1.1) circle (3.2pt);
\fill (6.0, 6.6) circle (3.2pt);
\fill (5.1, 2.0) circle (3.2pt);
\node[root] at (5.3, 2.9) {};

\end{tikzpicture}
\par\smallskip{\small A $k$-d tree of 15 points built by medians. The circled point is the root: its vertical line
splits the square in two. Each half is split by a horizontal line through its median point, and so on, alternating.
Blue: splits on $x$; red: splits on $y$. Drawn by \code{figs.py} from \code{kdtree.py}.}
\end{omfigure}

\begin{proposition}[Nearest-neighbour search in a $k$-d tree]\label{prop:b2:trees-and-clustering:kd-correct}
The function \code{nearest} of \cref{lst:b2:trees-and-clustering:kdtree} returns a point of the tree nearest to $q$.
\end{proposition}

\begin{proof}
It visits every node except the subtrees it skips, and keeps the closest point visited; it suffices that no skipped
subtree holds a point strictly closer than the best at the time of the skip, since the best only improves. Let the node
be $p$, splitting on axis $a$, and suppose $q_a \leq p_a$: the skipped subtree is the right one, whose points $r$ have
$r_a \geq p_a$, so $\lVert r - q \rVert^2 \geq (r_a - q_a)^2 \geq (p_a - q_a)^2$. The subtree is skipped only when $(p_a
- q_a)^2$ is at least the best squared distance, so no point of it is strictly closer. The case $q_a > p_a$ is the same
with the left subtree, whose points have $r_a \leq p_a < q_a$.
\end{proof}

\omoutput{code/bachelor-2/29-trees-and-clustering/outputs/kdtree.txt}{Nodes visited per query, against the $n$ distances
of the brute-force search, for uniform random points.}\label{out:b2:trees-and-clustering:kdtree}

On uniform random points in the plane, the search visits a few nodes per level of the tree: 23 out of 100,000 on
average. That is an observation, not a bound: points arranged to be all at nearly the same distance from the query force
the search to visit every node (\cref{exo:b2:trees-and-clustering:12}), and in many dimensions the pruning test rarely
succeeds, so $k$-d trees pay off only when the dimension is small.

\begin{history}
% ledger: b2:bentley-1975
Jon Bentley published multidimensional binary search trees, the $k$-d trees, in 1975.
\end{history}

\section{Decision trees}\label{sec:b2:trees-and-clustering:trees}

\begin{definition}[Decision tree]\label{def:b2:trees-and-clustering:decision-tree}
A \emph{decision tree}\index{decision tree} is a binary tree whose internal nodes carry a test $x_f \leq t$ on a \omterm{def:b2:trees-and-clustering:examples}{feature}
$f$ with a threshold $t$, and whose leaves carry a \omterm{def:b2:trees-and-clustering:examples}{class label}. It labels an \omterm{def:b2:trees-and-clustering:examples}{example} by following the branch of each
test, left when the test holds, down to a leaf.
\end{definition}

To choose the first question, a program needs to measure how mixed the labels are, and how much a question unmixes
them.

\begin{definition}[Entropy, information gain]\label{def:b2:trees-and-clustering:entropy}
The \emph{entropy}\index{entropy} of a finite probability distribution $(p_1, \dots, p_m)$ is $H = -\sum_i p_i \log_2
p_i$, with $0 \log_2 0 = 0$; the entropy of a set of labelled \omterm{def:b2:trees-and-clustering:examples}{examples} is that of the frequencies of its labels. A split
of a set $S$ into $S_1$ and $S_2$ has \emph{information gain}\index{information gain}
\[
G = H(S) - \frac{|S_1|}{|S|} H(S_1) - \frac{|S_2|}{|S|} H(S_2).
\]
\end{definition}

The \omterm{def:b2:trees-and-clustering:entropy}{entropy} is $0$ for a set with a single label and $\log_2 m$, its largest value, for $m$ labels in equal numbers: one
bit for a fair coin. The gain is how much the split lowers the \omterm{def:b2:trees-and-clustering:entropy}{entropy}, on average over the two parts.

\begin{omfigure}
\begin{tikzpicture}
\begin{axis}[width=6.5cm, height=4.6cm, xlabel={$p$, fraction of late loans}, ylabel={$h(p)$ in bits}, xmin=0, xmax=1,
  ymin=0, ymax=1.1, samples=101, domain=0.0001:0.9999, xtick={0, 0.25, 0.5, 0.75, 1}]
\addplot[omDef, thick] {-x*log2(x) - (1-x)*log2(1-x)};
\draw[omThm, thick] (axis cs:0.1, 0.469) -- (axis cs:0.7, 0.881);
\fill[omThm] (axis cs:0.1, 0.469) circle (1.5pt) (axis cs:0.7, 0.881) circle (1.5pt);
\fill[black] (axis cs:0.4, 0.971) circle (1.5pt);
\draw[densely dotted] (axis cs:0.4, 0.675) -- (axis cs:0.4, 0.971);
\end{axis}
\end{tikzpicture}
\par\smallskip{\small The \omterm{def:b2:trees-and-clustering:entropy}{entropy} of two labels, $h(p) = -p \log_2 p - (1 - p) \log_2(1 - p)$. A set with $40~\%$ late
loans split into halves with $10~\%$ and $70~\%$: the average \omterm{def:b2:trees-and-clustering:entropy}{entropy} of the parts, on the chord, lies below the
\omterm{def:b2:trees-and-clustering:entropy}{entropy} of the whole, on the curve. The dotted gap is the gain.}
\end{omfigure}

\begin{proposition}[Information gain is never negative]\label{prop:b2:trees-and-clustering:gain-nonneg}
Every split has $G \geq 0$.
\end{proposition}

\begin{proof}
For two labels: write $h(p) = -p \log_2 p - (1 - p) \log_2 (1 - p)$ for the \omterm{def:b2:trees-and-clustering:entropy}{entropy} of a set with a fraction $p$ of the
first label. Its second derivative on $(0, 1)$ is $-1/(p(1 - p) \ln 2) < 0$, so $h$ is concave on $[0, 1]$. If $S_1$ and
$S_2$ have fractions $p_1$ and $p_2$ and $w = |S_1|/|S|$, then $S$ has fraction $p = w p_1 + (1 - w) p_2$, and concavity
gives $h(p) \geq w\,h(p_1) + (1 - w)\,h(p_2)$, which is $G \geq 0$. With more labels the same argument applies to the
\omterm{def:b2:trees-and-clustering:entropy}{entropy} as a function of the vector of frequencies, which is concave too; we admit that case.
\end{proof}

\begin{method}[ID3 with binary splits]\label{met:b2:trees-and-clustering:id3}\index{ID3}
To grow a tree on a set $S$ of labelled \omterm{def:b2:trees-and-clustering:examples}{examples}: if no split of $S$ of the form $x_f \leq t$ has positive gain, or the
depth limit is reached, return a leaf with the majority label of $S$. Otherwise take the split of largest gain, with $t$
halfway between two consecutive values of $x_f$ in $S$, and grow the two subtrees on the two parts.
\end{method}

\omcode{code/bachelor-2/29-trees-and-clustering/listings/id3.py}{8}{31}{\omterm{def:b2:trees-and-clustering:entropy}{Entropy}, and the best split by one sweep over
each sorted \omterm{def:b2:trees-and-clustering:examples}{feature}}\label{lst:b2:trees-and-clustering:best-split}

\needspace{14\baselineskip}
\omcode{code/bachelor-2/29-trees-and-clustering/listings/id3.py}{34}{44}{ID3}\label{lst:b2:trees-and-clustering:id3}

Sorting the \omterm{def:b2:trees-and-clustering:examples}{examples} on each \omterm{def:b2:trees-and-clustering:examples}{feature} lets the sweep move one \omterm{def:b2:trees-and-clustering:examples}{example} at a time from right to left and update the label
counts in constant time, so the best split of $n$ \omterm{def:b2:trees-and-clustering:examples}{examples} on $d$ \omterm{def:b2:trees-and-clustering:examples}{features} costs $O(dn \log n)$. A split on a set
whose \omterm{def:b2:trees-and-clustering:examples}{examples} all carry the same label has no positive gain, so ID3 stops on pure sets; it also stops when no question
helps, which a depth limit makes earlier.

\begin{omfigure}
\begin{forest}
for tree={font=\small, draw, rounded corners, align=center, l sep=12pt, s sep=14pt, inner sep=3pt},
leaf/.style={draw=none, font=\small\itshape},
[{late returns $\leq 1.5$?}
  [{on time\\448 / 0}, leaf, edge label={node[pos=0.8, above left, inner sep=1pt, font=\scriptsize] {yes}}]
  [{distance $\leq 2490$?}, edge label={node[pos=0.8, above right, inner sep=1pt, font=\scriptsize] {no}}
    [{on time\\95 / 0}, leaf, edge label={node[pos=0.8, above left, inner sep=1pt, font=\scriptsize] {yes}}]
    [{late\\0 / 157}, leaf, edge label={node[pos=0.8, above right, inner sep=1pt, font=\scriptsize] {no}}]]]
\end{forest}
\par\smallskip{\small The tree of \cref{out:b2:trees-and-clustering:hook}, with the training loans that reach each leaf,
on time / late. The threshold 2490 lies halfway between the two nearest distances on either side of the planted 2500.}
\end{omfigure}

\begin{definition}[Overfitting]\label{def:b2:trees-and-clustering:overfitting}
A \omterm{def:b2:trees-and-clustering:learning}{classifier} \emph{overfits}\index{overfitting} its \omterm{def:b2:trees-and-clustering:examples}{training set} when its \omterm{def:b2:trees-and-clustering:confusion}{accuracy} there is much higher than on new
\omterm{def:b2:trees-and-clustering:examples}{examples} from the same source, and a simpler \omterm{def:b2:trees-and-clustering:learning}{classifier}, built from the same \omterm{def:b2:trees-and-clustering:examples}{training set}, does better on the new
\omterm{def:b2:trees-and-clustering:examples}{examples}.
\end{definition}

On noiseless loans, the tree stopped by itself at depth 2. On loans whose labels are flipped with probability $0.15$,
every flipped label is an \omterm{def:b2:cpp-classes-raii:exception}{exception} that a deeper tree can isolate with more questions:

\omcode{code/bachelor-2/29-trees-and-clustering/listings/overfit.py}{6}{11}{Trees of growing depth on noisy
loans}\label{lst:b2:trees-and-clustering:overfit}

\omoutput{code/bachelor-2/29-trees-and-clustering/outputs/overfit.txt}{Training \omterm{def:b2:trees-and-clustering:confusion}{accuracy} keeps rising; test \omterm{def:b2:trees-and-clustering:confusion}{accuracy} peaks
at depth 2.}\label{out:b2:trees-and-clustering:overfit}

\begin{omfigure}
\begin{tikzpicture}
\begin{axis}[width=0.7\linewidth, height=5.2cm, xlabel={depth limit of the tree}, ylabel={accuracy}, ymin=0.68,
  ymax=1.02, xmin=-0.5, xmax=16.5, legend pos=north west, legend style={font=\small}, xtick={0,2,4,6,8,10,12,16}]
\addplot[mark=*, omThm] table[x=depth, y=train] {code/bachelor-2/29-trees-and-clustering/figures/overfit.dat};
\addplot[mark=square*, omDef] table[x=depth, y=test] {code/bachelor-2/29-trees-and-clustering/figures/overfit.dat};
\legend{training set, test set}
\end{axis}
\end{tikzpicture}
\par\smallskip{\small \Cref{out:b2:trees-and-clustering:overfit} as a chart. Past depth 2, each extra level fits more of
the noise in the \omterm{def:b2:trees-and-clustering:examples}{training set}, and the tree does worse on new loans: at depth 16 it labels every training loan correctly
and is beaten by the rule ``always on time''.}
\end{omfigure}

The depth limit is not learned from the \omterm{def:b2:trees-and-clustering:examples}{training set}, which always prefers the deepest tree. It is chosen on \omterm{def:b2:trees-and-clustering:examples}{examples}
held out for that purpose, a validation set, or by cross-validation (\cref{exo:b2:trees-and-clustering:7}); the \omterm{def:b2:trees-and-clustering:examples}{test set}
is looked at once, at the end.

\begin{history}
% ledger: b2:quinlan-1986
Ross Quinlan published ID3, \omterm{def:b2:trees-and-clustering:decision-tree}{decision trees} grown by \omterm{def:b2:trees-and-clustering:entropy}{information gain}, in 1986.
\end{history}

\section{Clustering by $k$-means}\label{sec:b2:trees-and-clustering:kmeans}

\omimg{book4/ai/b2-trees-and-clustering-returns-desk.jpg}{0.6\linewidth}{A returns desk with its trolley. The loans
have labels that arrive by themselves, on time or late; the members who \omterm{def:b2:rust-ownership:borrow}{borrow} them do not come sorted into groups, and
the groups, if any, have to be found.}

\begin{definition}[Clustering, centroid, $k$-means cost]\label{def:b2:trees-and-clustering:clustering}
A \emph{clustering}\index{clustering} of a finite set of points $P \subset \mathbb{R}^d$ into $k$ \emph{clusters} is a
partition of $P$ into $k$ parts, given by an assignment $a : P \to \{1, \dots, k\}$. The \emph{centroid}\index{centroid}
of a non-empty finite set is the mean of its points. Given centres $c_1, \dots, c_k$, the \emph{$k$-means
cost}\index{k-means cost@$k$-means cost} of an assignment is $\sum_{p \in P} \lVert p - c_{a(p)} \rVert^2$.
\end{definition}

\begin{lemma}[The centroid minimises the squared distances]\label{lem:b2:trees-and-clustering:centroid}
For a non-empty finite set $C$ with \omterm{def:b2:trees-and-clustering:clustering}{centroid} $m$ and any point $c$, $\sum_{p \in C} \lVert p - c \rVert^2 = \sum_{p \in C}
\lVert p - m \rVert^2 + |C|\,\lVert m - c \rVert^2$. In particular $m$ is the unique point that minimises the sum.
\end{lemma}

\begin{proof}
Write $p - c = (p - m) + (m - c)$ and expand the square: the cross terms add up to $2 (m - c) \cdot \sum_{p}(p - m)$,
which is $0$ by the definition of $m$.
\end{proof}

\begin{method}[Lloyd's algorithm]\label{met:b2:trees-and-clustering:lloyd}\index{Lloyd's algorithm}
Start from $k$ centres, for \omterm{def:b2:trees-and-clustering:examples}{example} $k$ points of $P$ drawn at random. Repeat: assign each point to its nearest centre
(the lowest index on a tie); if the assignment has not changed, stop; otherwise move each centre to the \omterm{def:b2:trees-and-clustering:clustering}{centroid} of its
cluster, or leave it in place if its cluster is empty.
\end{method}

\omcode{code/bachelor-2/29-trees-and-clustering/listings/kmeans.py}{12}{39}{\omterm{met:b2:trees-and-clustering:lloyd}{Lloyd's algorithm}, recording the cost after
each assignment}\label{lst:b2:trees-and-clustering:lloyd}

\begin{theorem}[Lloyd's algorithm stops]\label{thm:b2:trees-and-clustering:lloyd-stops}
The cost never increases from one assignment to the next, and the algorithm stops after finitely many iterations.
\end{theorem}

\begin{proof}
\Cref{exo:b2:trees-and-clustering:8}: the assignment step can only lower each point's term, the \omterm{def:b2:trees-and-clustering:clustering}{centroid} step can only
lower each cluster's sum by \cref{lem:b2:trees-and-clustering:centroid}, and an assignment never comes back.
\end{proof}

It stops, but not necessarily at the best \omterm{def:b2:trees-and-clustering:clustering}{clustering}: the result depends on the start.

\omoutput{code/bachelor-2/29-trees-and-clustering/outputs/kmeans.txt}{One run, then the final costs of 20 runs from
different random starts on the same 150 points.}\label{out:b2:trees-and-clustering:kmeans}

\begin{omfigure}
\begin{tikzpicture}[scale=0.8,
  base/.style={inner sep=0pt, minimum size=2.6pt},
  pt0/.style={base, circle, fill=omDef}, pt1/.style={base, rectangle, fill=omThm},
  pt2/.style={base, diamond, fill=omProp, minimum size=3.4pt}, pt3/.style={base, circle, fill=omMeth},
  pt4/.style={base, isosceles triangle, isosceles triangle apex angle=60, shape border rotate=90, fill=black!55,
    minimum size=3.4pt},
  centre/.style={draw=black, very thick, cross out, minimum size=7pt, inner sep=0pt}]
\foreach \k/\t in {0/{start}, 1/{after one step}, 2/{final, cost 625.7}} {
  \begin{scope}[xshift=\k*6.4cm]
  \draw[black!40] (-0.7, 0.2) rectangle (5.5, 5.4);
  \input{code/bachelor-2/29-trees-and-clustering/figures/kmeans_\k.tex}
  \node[font=\small] at (2.4, -0.2) {\t};
  \end{scope}}
\end{tikzpicture}
\par\smallskip{\small \omterm{met:b2:trees-and-clustering:lloyd}{Lloyd's algorithm} from the start of seed 3, on five planted groups of 30 points. Crosses: the
centres; colours and shapes: the clusters. Two centres start in the upper right group and never leave it, and one
centre has to cover the two lower groups: a local minimum at cost 625.7, against 131.2 for the planted grouping. Drawn by \code{figs.py}.}
\end{omfigure}

Eight starts of twenty find the planted groups; the others stop at costs up to five times higher. The usual remedy is to
run the algorithm from many starts and keep the best, and to start better. The \emph{$k$-means++}\index{k-means++@$k$-means++}
seeding chooses the first centre at random and each next one among the points with probability proportional to its
squared distance to the nearest centre chosen so far, so that far groups are likely to get a centre:

\omcode{code/bachelor-2/29-trees-and-clustering/listings/kmeanspp.py}{8}{15}{The $k$-means++ seeding}\label{lst:b2:trees-and-clustering:kmeanspp}

\omoutput{code/bachelor-2/29-trees-and-clustering/outputs/kmeanspp.txt}{Twenty runs from each seeding: how often each
finds the planted groups, and the costs.}\label{out:b2:trees-and-clustering:kmeanspp}

Better starts also need fewer steps. The seeding does not remove local minima, only makes them rarer: over 100 seeds,
\cref{exo:b2:trees-and-clustering:11} counts them.

\begin{history}
% ledger: b2:lloyd-1957
% ledger: b2:lloyd-1982
Stuart Lloyd proposed the algorithm in 1957, for the quantisation of signals; his work circulated widely but was
published only in 1982.
\end{history}

\section{Hierarchical clustering}\label{sec:b2:trees-and-clustering:hac}

$k$-means asks for $k$ in advance. Hierarchical \omterm{def:b2:trees-and-clustering:clustering}{clustering} builds every number of clusters at once.

\begin{definition}[Hierarchical agglomerative clustering, linkage, dendrogram]\label{def:b2:trees-and-clustering:hac}
\emph{Hierarchical agglomerative clustering}\index{hierarchical agglomerative clustering} starts with each point in a
cluster of its own and repeatedly merges the two clusters at the least \emph{linkage}\index{linkage} distance, until one
cluster is left. The \emph{single}, \emph{complete} and \emph{average} linkages of two clusters are the least, the
largest and the mean distance between a point of one and a point of the other. The merges, drawn as a tree whose leaves
are the points and whose internal nodes sit at the height of their merge, form the
\emph{dendrogram}\index{dendrogram}; cutting it after $n - k$ merges gives $k$ clusters.
\end{definition}

\omcode{code/bachelor-2/29-trees-and-clustering/listings/hac.py}{14}{26}{Agglomerative \omterm{def:b2:trees-and-clustering:clustering}{clustering} with any \omterm{def:b2:trees-and-clustering:hac}{linkage}, by
brute force}\label{lst:b2:trees-and-clustering:hac}

\omoutput{code/bachelor-2/29-trees-and-clustering/outputs/hac.txt}{Twelve points, three \omterm{def:b2:trees-and-clustering:hac}{linkages}, and single \omterm{def:b2:trees-and-clustering:hac}{linkage}
against the spanning tree.}\label{out:b2:trees-and-clustering:hac}

\begin{omfigure}
\begin{tikzpicture}[x=0.42cm, y=0.42cm]
\draw[black!50] (-0.6, -0.6) rectangle (9.6, 9.6);
\foreach \x/\y/\n in {1.0/1.0/A, 2.2/1.4/B, 1.3/2.9/C, 4.6/1.6/D, 6.8/1.1/E, 8.1/2.2/F, 9.0/0.9/G, 7.9/3.9/H,
  4.2/6.8/I, 5.1/8.9/J, 6.3/7.6/K, 3.0/8.7/L} {
  \fill (\x, \y) circle (2pt); \node[font=\scriptsize, above right=-1pt] at (\x, \y) {\n};}
\begin{scope}[xshift=5.4cm, x=0.5cm, y=1.1cm, leaf/.style={font=\small, below}]
% written by figs.py; do not edit
\node[leaf] at (0, 0) {I};
\node[leaf] at (1, 0) {L};
\node[leaf] at (2, 0) {J};
\node[leaf] at (3, 0) {K};
\node[leaf] at (4, 0) {C};
\node[leaf] at (5, 0) {A};
\node[leaf] at (6, 0) {B};
\node[leaf] at (7, 0) {D};
\node[leaf] at (8, 0) {H};
\node[leaf] at (9, 0) {E};
\node[leaf] at (10, 0) {F};
\node[leaf] at (11, 0) {G};
\draw (2, 0.00) |- (3, 1.77) -- (3, 0.00);
\draw (1, 0.00) |- (2.5, 2.11) -- (2.5, 1.77);
\draw (0, 0.00) |- (1.75, 2.25) -- (1.75, 2.11);
\draw (5, 0.00) |- (6, 1.26) -- (6, 0.00);
\draw (4, 0.00) |- (5.5, 1.75) -- (5.5, 1.26);
\draw (10, 0.00) |- (11, 1.58) -- (11, 0.00);
\draw (9, 0.00) |- (10.5, 1.70) -- (10.5, 1.58);
\draw (8, 0.00) |- (9.75, 1.71) -- (9.75, 1.70);
\draw (7, 0.00) |- (8.875, 2.26) -- (8.875, 1.71);
\draw (4.75, 1.75) |- (7.9375, 2.41) -- (7.9375, 2.26);
\draw (0.875, 2.25) |- (6.34375, 4.03) -- (6.34375, 2.41);

\draw[->] (-1, 0) -- (-1, 4.4) node[above, font=\scriptsize] {height};
\foreach \h in {0, 1, 2, 3, 4} \draw (-1.1, \h) -- (-0.9, \h) node[left=2pt, font=\scriptsize] {\h};
\draw[omThm, dashed] (-0.5, 2.33) -- (11.5, 2.33) node[right, font=\scriptsize] {3 clusters};
\end{scope}
\end{tikzpicture}
\par\smallskip{\small Left: the twelve points. Right: their \omterm{def:b2:trees-and-clustering:hac}{single-linkage} \omterm{def:b2:trees-and-clustering:hac}{dendrogram}, drawn by \code{figs.py} from
\code{hac.py}. Cutting between the heights 2.26 and 2.41 leaves three clusters: D joins E, F, G and H through its
nearest neighbour, at 2.26, before A, B, C reach them at 2.41. Complete \omterm{def:b2:trees-and-clustering:hac}{linkage} puts D with A, B and C instead.}
\end{omfigure}

Single \omterm{def:b2:trees-and-clustering:hac}{linkage} follows chains of close points, which finds groups of any shape but may join two groups through a few
points between them; complete \omterm{def:b2:trees-and-clustering:hac}{linkage} prefers compact groups of small diameter. The brute-force listing recomputes
every \omterm{def:b2:trees-and-clustering:hac}{linkage} at every merge: each of the $n - 1$ merges looks at every pair of points in different clusters, $O(n^3)$
distances in all, which is fine for a dozen points only. Single \omterm{def:b2:trees-and-clustering:hac}{linkage} needs no such work: it is already known from the Year 1 volume.

\begin{theorem}[Single linkage and the minimum spanning tree]\label{thm:b2:trees-and-clustering:single-linkage}
Let the $\binom{n}{2}$ distances between $n$ points be distinct. For each $k$, the $k$ clusters of single \omterm{def:b2:trees-and-clustering:hac}{linkage} are the
connected components of the minimum spanning tree of the complete graph on the points, weighted by distance, once its
$k - 1$ heaviest edges are removed; equivalently, the components after Kruskal's algorithm has accepted $n - k$ edges.
\end{theorem}

\begin{proof}
\Cref{exo:b2:trees-and-clustering:10}. The idea: the single \omterm{def:b2:trees-and-clustering:hac}{linkage} of two clusters is the shortest edge between them, and
the shortest edge between two different clusters is exactly the next edge that Kruskal's algorithm accepts.
\end{proof}

So single \omterm{def:b2:trees-and-clustering:hac}{linkage} costs a sort of the $\binom{n}{2}$ distances and a union--find, $O(n^2 \log n)$ in all, and its
clusters for every $k$ come out of one spanning tree.

\section{Exercises}\label{sec:b2:trees-and-clustering:exercises}

\begin{exercise}[$\star$]\label{exo:b2:trees-and-clustering:1}
From the \omterm{def:b2:trees-and-clustering:confusion}{confusion matrix} of \cref{out:b2:trees-and-clustering:confusion}, compute the \omterm{def:b2:trees-and-clustering:confusion}{accuracy}, the fraction of late
loans that are caught, and the fraction of ``late'' predictions that are right. Which of the three should the library
watch if a false alarm costs a polite phone call and a missed late loan costs nothing?
\end{exercise}

\begin{exercise}[$\star$]\label{exo:b2:trees-and-clustering:2}
Run the $k$-nearest-neighbour \omterm{def:b2:trees-and-clustering:learning}{classifier} on the raw \omterm{def:b2:trees-and-clustering:examples}{features}, without scaling, with \code{exo_scaling.py}. Which
\omterm{def:b2:trees-and-clustering:examples}{feature} decides the distance, and why does \omterm{def:b2:trees-and-clustering:confusion}{accuracy} drop?
\end{exercise}

\begin{exercise}[$\star$]\label{exo:b2:trees-and-clustering:3}
Write a search for all the points of a $k$-d tree inside a rectangle, which skips the subtrees that cannot meet it.
Compare the nodes visited with the points reported, using \code{exo_range.py}.
\end{exercise}

\begin{exercise}[$\star$]\label{exo:b2:trees-and-clustering:4}
A set of 20 loans has 8 late ones. A split sends 10 loans to the left, 7 of them late, and 10 to the right, 1 of them
late. Compute its \omterm{def:b2:trees-and-clustering:entropy}{information gain} by hand, to three decimals.
\end{exercise}

\begin{exercise}[$\star\star$]\label{exo:b2:trees-and-clustering:5}
Run one iteration of \omterm{met:b2:trees-and-clustering:lloyd}{Lloyd's algorithm} by hand on the twelve points of \cref{out:b2:trees-and-clustering:hac}, from the
centres A, E and I: the assignment, the new centres and the two costs. Check with \code{exo_lloyd_step.py}.
\end{exercise}

\begin{exercise}[$\star\star$]\label{exo:b2:trees-and-clustering:6}
Some versions of ID3 split a \omterm{def:b2:trees-and-clustering:examples}{feature} with one branch per value. Add the loan's number, different for each loan, as a
\omterm{def:b2:trees-and-clustering:examples}{feature}, and compute the gain of such a split on it, on the days lent and on the late returns. What goes wrong? Is the
\emph{gain ratio}, the gain divided by the \omterm{def:b2:trees-and-clustering:entropy}{entropy} of the branch sizes, enough to fix it here? Use
\code{exo_identifier.py}.
\end{exercise}

\begin{exercise}[$\star\star$]\label{exo:b2:trees-and-clustering:7}
In \emph{5-fold cross-validation}, the \omterm{def:b2:trees-and-clustering:examples}{training set} is cut into five parts; each part in turn serves as the validation
set for a tree grown on the other four, and the five accuracies are averaged. Use it to choose the depth of the tree on
the 700 noisy training loans, with \code{exo_crossval.py}. Compare with \cref{out:b2:trees-and-clustering:overfit}.
\end{exercise}

\begin{exercise}[$\star\star$]\label{exo:b2:trees-and-clustering:8}
Prove \cref{thm:b2:trees-and-clustering:lloyd-stops}.
\end{exercise}

\begin{exercise}[$\star\star$]\label{exo:b2:trees-and-clustering:9}
Run complete \omterm{def:b2:trees-and-clustering:hac}{linkage} by hand on the twelve points for the first five merges, then compare the full list of merges with
\code{exo_complete.py}. Why does D end with A, B and C, while single \omterm{def:b2:trees-and-clustering:hac}{linkage} put it with E, F, G and H?
\end{exercise}

\begin{exercise}[$\star\star\star$]\label{exo:b2:trees-and-clustering:10}
Prove \cref{thm:b2:trees-and-clustering:single-linkage}.
\end{exercise}

\begin{exercise}[$\star\star\star$]\label{exo:b2:trees-and-clustering:11}
Compare uniform and $k$-means++ seeding over 100 seeds each with \code{exo_kmeanspp.py}: how often does each find the
best \omterm{def:b2:trees-and-clustering:clustering}{clustering}, and how many different local minima appear?
\end{exercise}

\begin{exercise}[$\star\star\star$]\label{exo:b2:trees-and-clustering:12}
Build a set of points and a query for which the nearest-neighbour search of \cref{lst:b2:trees-and-clustering:kdtree}
visits every node, and explain why. Check with \code{exo_kd_worst.py}.
\end{exercise}

\section{Problem: Which loans come back late?}\label{sec:b2:trees-and-clustering:problem}

\begin{problem}[{Weekend problem --- Which loans come back late?}]\label{pb:b2:trees-and-clustering:1}
The library has 1,000 past loans, with the same five \omterm{def:b2:trees-and-clustering:examples}{features} and a label that is wrong one time in ten. It wants a
tree it can trust, and a first look at its members' habits.

\omcode{code/bachelor-2/29-trees-and-clustering/listings/loans_problem.py}{9}{27}{The split, the depth chosen on
validation, and the test}\label{lst:b2:trees-and-clustering:problem}

\omoutput{code/bachelor-2/29-trees-and-clustering/outputs/loans_problem.txt}{The data, the tree, its test, and the
members clustered for $k = 1$ to $6$.}\label{out:b2:trees-and-clustering:problem}

\medskip\noindent\textbf{Part I --- The data and the split.}\par\nopagebreak
\begin{enumerate}
\item Why are the loans shuffled before they are split?
\item What is each of the three parts used for, and why must the \omterm{def:b2:trees-and-clustering:examples}{test set} be used only once?
\item The \omterm{def:b2:trees-and-clustering:examples}{test set} has 21.7~\% late loans, the \omterm{def:b2:trees-and-clustering:examples}{training set} 30.2~\%. Is that a problem? What would a split that keeps the
same proportion in every part do?
\item With labels wrong one time in ten, what is the best \omterm{def:b2:trees-and-clustering:confusion}{accuracy} any \omterm{def:b2:trees-and-clustering:learning}{classifier} can expect on the \omterm{def:b2:trees-and-clustering:examples}{test set}?
\item What is the majority baseline on the \omterm{def:b2:trees-and-clustering:examples}{test set}, and why is it the first number to compare with?
\item Which of the five \omterm{def:b2:trees-and-clustering:examples}{features} would you expect to matter, before looking at the tree?
\end{enumerate}

\medskip\noindent\textbf{Part II --- A tree that finds the planted rule.}\par\nopagebreak
\begin{enumerate}[resume, beginpenalty=10000]
\item Which depth does validation choose, and why the shallowest among the best?
\item Read the tree: which loans does it call late?
\item Compare it with the planted rule of \code{dataset.py}. Why does its agreement with the rule reach 1.000?
\item The left branch splits on the pages, yet both its leaves say ``on time''. Why did ID3 make that split, and what
does it change?
\item Compute the \omterm{def:b2:trees-and-clustering:entropy}{information gain} of the root split from the counts in the leaves.
\item State the problem's result: test \omterm{def:b2:trees-and-clustering:confusion}{accuracy} against the majority baseline.
\end{enumerate}

\medskip\noindent\textbf{Part III --- Overfitting with noise.}\par\nopagebreak
\begin{enumerate}[resume, beginpenalty=10000]
\item Validation \omterm{def:b2:trees-and-clustering:confusion}{accuracy} falls from depth 4 on. Explain it with the noise.
\item Why is the tree finally grown on the training and validation sets together?
\item The test \omterm{def:b2:trees-and-clustering:confusion}{accuracy}, 0.923, is above the 0.9 of question 4. Is that possible? What does it say about a single \omterm{def:b2:trees-and-clustering:examples}{test
set} of 300?
\item Would $k$-nearest neighbours with all five \omterm{def:b2:trees-and-clustering:examples}{features} do as well? Use \cref{out:b2:trees-and-clustering:knn}.
\item How would the tree change if the noise were 30~\%?
\item Propose a check that the tree's rule is stable: what would you rerun, and what would you compare?
\end{enumerate}

\medskip\noindent\textbf{Part IV --- The members, clustered.}\par\nopagebreak
\begin{enumerate}[resume, beginpenalty=10000]
\item Each member is a point (loans a year, mean days kept). Why is \omterm{def:b2:trees-and-clustering:clustering}{clustering} here unsupervised?
\item The cost drops from 10,797 to 2,701 to 810, then slowly. Which $k$ would you choose, and why?
\item Describe the three groups from their centres.
\item Why does the program keep the best of 10 starts for each $k$?
\item Would the two \omterm{def:b2:trees-and-clustering:examples}{features} need scaling here? Compare their ranges.
\item Which other method of the chapter could find the same groups without choosing $k$ first?
\end{enumerate}
\end{problem}
